\documentclass[11pt,letterpaper]{article}

\usepackage[spanish,es-noshorthands,es-tabla]{babel}
\usepackage{fontspec}
\setmainfont{Source Serif Pro}
\usepackage{microtype}
\usepackage{geometry}
\usepackage{enumitem}
\usepackage{xcolor}
\usepackage{csquotes}
\usepackage{xurl}
\usepackage[hidelinks]{hyperref}
\usepackage[backend=biber,style=apa,sorting=nyt,maxcitenames=2]{biblatex}
\DeclareLanguageMapping{spanish}{spanish-apa}

\geometry{margin=2.3cm}
\hbadness=10000
\setlength{\emergencystretch}{1em}
\setlength{\parindent}{0pt}
\setlength{\parskip}{0.5em}
\setlist{nosep}
\addbibresource{referencias.bib}
\AtBeginBibliography{\sloppy\small}
\setlength{\bibitemsep}{0.3em}

\definecolor{cEje}{RGB}{26,58,102}
\definecolor{cGris}{RGB}{120,116,110}

% Enunciado de cada pregunta, en gris y con sangría.
\newcommand{\pregunta}[1]{%
  \begin{quote}\color{cGris}\small\itshape #1\end{quote}}

\begin{document}

% --- Portada -----------------------------------------------------------------
\begin{titlepage}
\centering
\vspace*{3.2cm}
{\large\scshape Inteligencia Artificial Generativa\\
Aplicada a la Transformación Digital\par}
\vspace{0.5cm}
{\normalsize Módulo 1. Actividad 1.1\par}
\vspace{2.6cm}
{\color{cEje}\rule{0.55\textwidth}{0.8pt}\par}
\vspace{0.7cm}
{\LARGE\bfseries La IA al alcance de todos\par}
\vspace{0.5cm}
{\large Accesibilidad, liderazgo en T\\
y lo que espero de este curso\par}
\vspace{0.7cm}
{\color{cEje}\rule{0.55\textwidth}{0.8pt}\par}
\vfill
{\large Carlos Luis Rojas Aragonés\par}
\vspace{0.3cm}
{\normalsize Chief Technology Officer Program\par}
\vspace{0.3cm}
{\normalsize 3 de noviembre de 2026\par}
\vspace*{1.5cm}
\end{titlepage}

% --- Pregunta 1 --------------------------------------------------------------
\section*{1. Una tecnología que se volvió accesible para todos}

\pregunta{Susan Doniz sostiene que el verdadero avance de la IA generativa fue
hacer que una tecnología muy potente resultara accesible para todo el mundo y
no únicamente para especialistas. Describa una situación en la que una
tecnología se volviera realmente accesible de este modo en su sector. ¿Qué
cambió para su organización o para su sector? ¿Qué ventajas aportó desde el
punto de vista del negocio?}

Lo contesto con un ejemplo: ChatGPT. El hecho de que llegara a tantos millones
de usuarios en tiempo récord no fue casualidad. Según un análisis de UBS
citado por Reuters, alcanzó unos 100 millones de usuarios mensuales en apenas
dos meses; a TikTok le tomó nueve meses llegar a esa cifra y a Instagram unos
dos años y medio \parencite{hu2023}. Para que esto se diera tuvieron que
coincidir varias cosas:

\begin{enumerate}
  \item \textbf{Es una tecnología diseñada para una persona de cualquier
    edad.} Si alguien sabe usar WhatsApp, entonces sabe usar ChatGPT. La
    interfaz es una conversación: no hay que aprender un lenguaje de
    programación ni un menú nuevo. Antes, usar un modelo de lenguaje
    requería un equipo de ciencia de datos; ahora basta con escribir.
  \item \textbf{Se aplica en cualquier negocio del mundo y brinda una ventaja
    competitiva increíble.} En un estudio con más de 5000 agentes de servicio
    al cliente, un asistente de IA generativa aumentó en promedio un 15\,\% los
    casos resueltos por hora, y en un 34\,\% los de los agentes con menos
    experiencia \parencite{brynjolfsson2025}. Es decir, la herramienta
    transfiere a los nuevos parte del conocimiento de los mejores.
  \item \textbf{Fue una tecnología verdaderamente disruptiva, que sorprendió
    tanto a usuarios no técnicos como a técnicos.} En un experimento con 758
    consultores de BCG, quienes usaron GPT-4 completaron un 12,2\,\% más de
    tareas, un 25,1\,\% más rápido y con una calidad más de un 40\,\% superior
    \parencite{dellacqua2023}. Los mismos autores advierten que, fuera de las
    tareas para las que la IA es buena, los consultores fallaron más; por eso
    hablan de una \enquote{frontera irregular}.
\end{enumerate}

¿Qué cambió para el sector? Que la IA dejó de ser un proyecto del área de
datos y pasó a estar en manos de todos. McKinsey reporta que en 2025 el 88\,\%
de las organizaciones ya usaba IA de forma regular en al menos una función del
negocio, frente al 78\,\% del año anterior \parencite{mckinsey2025}. La ventaja
para el negocio no está solamente en hacer lo mismo más rápido, sino en que
personas sin perfil técnico pueden probar ideas, automatizar tareas y
prototipar sin depender de un equipo de desarrollo. El mismo informe aclara
que la mayoría de las empresas todavía está en etapa de pilotos, así que el
reto ahora no es el acceso, sino escalar.

% --- Pregunta 2 --------------------------------------------------------------
\section*{2. Mi preparación para la IA según el liderazgo en T}

\pregunta{Susan Doniz presenta el modelo del liderazgo en T, que combina
profundidad técnica y una amplia visión de negocio. Evalúe su nivel actual de
preparación para la IA utilizando este modelo. ¿Cuáles considera que son sus
principales fortalezas? ¿En qué aspectos identifica oportunidades de mejora?
¿Qué supondría, en su caso, reforzar la dimensión menos desarrollada de la T?}

Mi principal fortaleza es que llevo más de seis años preparándome en
inteligencia artificial enfocada 100\,\% en negocios. Eso me ha permitido ver
muchos casos de éxito de la vida real, y también casos no exitosos. Los
segundos enseñan tanto como los primeros: según un estudio de MIT Sloan
Management Review y BCG, solo un 10\,\% de las empresas obtenía beneficios
financieros significativos con IA, y lo que las distinguía no era la
tecnología, sino que cambiaban sus procesos para aprender junto con ella
\parencite{ransbotham2020}.

El modelo en T me llama la atención porque se parece a algo que una vez me
dijo un amigo: \enquote{Un CTO es una persona con un mar de conocimiento, pero
de un centímetro de profundidad}. La idea tampoco es nueva en la gestión:
Hansen y von Oetinger describieron hace más de veinte años al gerente en T como
alguien comprometido a fondo con su área (la parte vertical) que además
comparte conocimiento con el resto de la organización (la parte horizontal)
\parencite{hansen2001}.

Creo que debemos elegir muy bien en qué temas profundizar y en cuáles no. Eso
nos convierte en especialistas de algunas áreas y nos deja un conocimiento
general de otras. Sin perder nunca de vista la forma en que funcionan los
negocios y cómo hacen dinero, porque eso marca una gran diferencia a la hora
de presentar un proyecto de tecnología ante una junta directiva.

La mejor analogía que puedo dar es la de los médicos generales y los
especialistas. El médico general no sabe operar un corazón, pero sabe
reconocer cuándo hace falta un cardiólogo y a cuál enviar al paciente. David
Epstein defiende algo parecido: en entornos cambiantes, donde las reglas no
están claras, quienes tienen amplitud de experiencias suelen adaptarse mejor
que quienes solo han profundizado en una cosa \parencite{epstein2019}. La IA
es justamente uno de esos entornos.

En mi caso, la parte horizontal de la T (negocio) es la más desarrollada. La
oportunidad de mejora está en mantener al día la parte vertical: la IA
generativa cambia cada pocos meses, y reforzarla significa reservar tiempo
para probar los modelos y herramientas nuevas con las manos, no solo leer
sobre ellas, de forma que pueda evaluar con criterio propio qué es útil y qué
es ruido.

% --- Pregunta 3 --------------------------------------------------------------
\section*{3. Lo que quiero lograr al terminar el curso}

\pregunta{Al finalizar este curso, ¿qué le gustaría ser capaz de hacer, decidir
o liderar que hoy todavía no puede? Sea específico e identifique un reto, una
decisión o una conversación real de su organización para la que este curso
pueda ayudarle a prepararse.}

Me gustaría poder participar en cualquier foro de IA del mundo y que mi aporte
sea significativo para la audiencia.

Para mí, un aporte significativo es uno que va más allá de repetir lo que ya
dicen las noticias. Hoy casi todas las empresas usan IA \parencite{mckinsey2025},
pero muy pocas logran resultados financieros con ella \parencite{ransbotham2020}.
La conversación que vale la pena en esos foros es justamente esa: por qué
unas empresas sí y otras no, y qué decisiones de liderazgo marcan la
diferencia. Espero que este curso me dé el marco y los casos para hablar de
eso con datos y con criterio propio.

\printbibliography[title={Referencias}]

\end{document}
